\chapter{Heart Transplant}

\section*{Definition and Overview}
A heart transplant is an operation in which a diseased, failing heart is replaced with a healthy heart from a deceased organ donor. It is the definitive treatment for end-stage heart failure when all other therapies, including medication and mechanical support devices, can no longer maintain an acceptable quality of life or survival. Heart transplantation is a major operation followed by lifelong immunosuppression to prevent rejection. With modern care, most recipients survive for many years and return to a full, active life.

\section*{Epidemiology}
Heart transplantation is a relatively rare procedure: around 80--100 heart transplants are performed in Australia each year, at specialist transplant centres. Demand substantially exceeds the supply of donor organs, and patients may wait months to years on the transplant list. Around half of adults on the Australian waiting list receive a transplant, with others dying or being removed from the list while waiting. Survival after transplantation is excellent for a procedure of this magnitude: approximately 90\% of recipients survive at least one year and around 70--75\% survive at least ten years.

\section*{Aetiology and Pathophysiology}
\begin{itemize}
    \item \textbf{Indications:} end-stage heart failure due to cardiomyopathy, coronary artery disease, congenital heart disease, or valvular disease that has failed conventional therapy
    \item \textbf{Contraindications:} active infection or cancer, severe irreversible disease of other organs, significant pulmonary hypertension, and inability to adhere to the demanding immunosuppression regimen
    \item \textbf{Referral:} patients are referred when estimated survival without transplant is poor despite optimal medical and device therapy
    \item \textbf{Donor matching:} compatibility of blood type and size is required, with cross-matching for antibodies
    \item \textbf{The procedure:} the failing heart is removed and the donor heart is connected to the great vessels; most recipients remain on cardiopulmonary bypass during the transplant
\end{itemize}

\section*{Clinical Features}
\begin{itemize}
    \item Candidates have severe, persistent symptoms of heart failure: breathlessness at rest, fatigue, fluid retention, and reduced exercise tolerance
    \item Dependence on inotropes or mechanical support such as a ventricular assist device may be present
    \item Repeated hospital admissions for decompensated heart failure
    \item After transplant, symptoms of heart failure resolve as the new heart functions normally
    \item The transplanted heart is denervated: it beats at a higher baseline rate and responds to exercise more slowly
\end{itemize}

\section*{Differential Diagnosis}
\begin{itemize}
    \item Not every patient with advanced heart failure is a transplant candidate, and other options are always considered first
    \item Mechanical circulatory support, particularly left ventricular assist devices, may be used as a bridge to transplant or as long-term therapy
    \item Optimal medical therapy with modern heart failure drugs can stabilise some patients without transplant
    \item Palliative care is appropriate for patients not suitable for transplant or support devices
    \item The decision is made by a specialist transplant team after comprehensive assessment
\end{itemize}

\section*{When to Seek Medical Care}
Patients on the transplant waiting list are under close specialist supervision and must report any new infection, fever, or deterioration. After transplant, recipients must seek urgent review for fever, breathlessness, or any sign of rejection or infection.

\textbf{Call 000 (or local emergency services) for:}
\begin{itemize}
    \item Sudden breathlessness, chest pain, or fainting after transplant
    \item High fever, which may indicate infection in an immunosuppressed patient
    \item Signs of stroke or sudden neurological symptoms
    \item Heavy bleeding or other emergencies in the early postoperative period
\end{itemize}

\section*{Diagnosis and Investigations}
\begin{itemize}
    \item \textbf{Comprehensive transplant assessment:} cardiac, pulmonary, renal, liver, and psychosocial evaluation
    \item \textbf{Imaging and catheter studies:} right heart catheterisation to measure pulmonary pressures, plus echocardiography
    \item \textbf{Blood tests:} organ function, infection screening, blood type, and antibody levels
    \item \textbf{After transplant:} surveillance is vital, with endomyocardial biopsy used to detect rejection in the first year
    \item \textbf{Coronary angiography:} monitors for cardiac allograft vasculopathy, a form of accelerated coronary disease in the transplanted heart
    \item \textbf{Monitoring of immunosuppression:} drug levels are measured to keep rejection and toxicity in balance
\end{itemize}

\section*{Management}
\begin{itemize}
    \item \textbf{The transplant operation:} performed at a specialist centre with experienced surgical and intensive care teams
    \item \textbf{Immunosuppression:} a lifelong combination of medications to prevent rejection of the donor heart
    \item \textbf{Infection prevention:} vaccination, infection surveillance, and prompt treatment of any infection
    \item \textbf{Complication surveillance:} regular biopsy, echocardiography, and coronary assessment
    \item \textbf{Rehabilitation:} supervised exercise and cardiac rehabilitation to restore fitness
    \item \textbf{Cardiovascular risk management:} control of blood pressure, cholesterol, and diabetes, which are promoted by immunosuppressants
    \item \textbf{Long-term follow-up:} lifelong specialist care with coordinated transplant and general health services
\end{itemize}

\section*{Complications}
\begin{itemize}
    \item Rejection of the donor heart, most common in the first year and treated with intensified immunosuppression
    \item Infection due to immunosuppression, including serious bacterial, viral, and fungal infections
    \item Cardiac allograft vasculopathy, the main long-term limitation to survival
    \item Kidney disease caused by immunosuppressive drugs
    \item Malignancy, including skin cancer and lymphoproliferative disorders, from long-term immunosuppression
    \item Diabetes, hypertension, and weight gain as side effects of medications
    \item Psychological challenges, including anxiety, depression, and adjustment difficulties
\end{itemize}

\section*{Prognosis and Outcomes}
Heart transplantation offers the best survival and quality of life for eligible patients with end-stage heart failure. Median survival is now more than 12 years, and many recipients live for 20 years or more, returning to work, exercise, and family life. The first year is the highest risk period, dominated by rejection and infection, after which the risk falls substantially. Long-term outcomes are shaped by adherence to immunosuppression, management of side effects, and surveillance for cardiac allograft vasculopathy. Quality of life after the first year is typically excellent.

\section*{Prevention and Self-Care}
\begin{itemize}
    \item Register as an organ donor and discuss donation wishes with family
    \item For recipients: take immunosuppressive medication exactly as prescribed, without missing doses
    \item Attend all scheduled surveillance tests and clinic visits
    \item Practise scrupulous hygiene and report any fever or infection promptly
    \item Protect skin from sun exposure and have regular skin checks
    \item Maintain a heart-healthy lifestyle with diet, activity, and no smoking
    \item Seek psychological support to manage the emotional demands of transplantation
\end{itemize}

\section*{Sources and Further Reading}
The following reputable sources provide further information for healthcare students and patients seeking reliable, current advice about heart transplantation.

\begin{itemize}
    \item DonateLife. \textit{Heart transplantation and organ donation.} https://www.donatelife.gov.au/
    \item Heart Foundation. \textit{Heart transplant: what to expect.} https://www.heartfoundation.org.au/
    \item Healthdirect. \textit{Heart transplant.} https://www.healthdirect.gov.au/heart-transplant
    \item NHS. \textit{Heart transplant: overview and recovery.} https://www.nhs.uk/conditions/heart-transplant/
    \item Mayo Clinic. \textit{Heart transplant: why it's done and what to expect.} https://www.mayoclinic.org/
    \item MedlinePlus. \textit{Heart transplantation.} https://medlineplus.gov/hearttransplantation.html
\end{itemize}
